% !PRINTED 547
\weberpage{547}

If we apply the theorem of \S.\,164, 3. to only two factors, it follows that
% [sic] authority prints the Latin u as the exponent; the surrounding formula requires \mu.
\[
(\varepsilon,\alpha)^{\nu}(\varepsilon^{\lambda},\alpha)^{u}
\]
is a function in $\Omega(\varepsilon)$ when $\nu + \lambda\mu \equiv 0 \pmod{m}$. If we therefore set $\mu = 1$, $\nu = m - \lambda$, it follows, if
\[
\chi_\lambda = b_0^{(\lambda)} + b_1^{(\lambda)}\varepsilon + \cdots + b_{m-1}^{(\lambda)}\varepsilon^{m-1}
\]
denotes a quantity in $\Omega(\varepsilon)$, that
\[
(\varepsilon,\alpha)^{m-\lambda}(\varepsilon^{\lambda},\alpha) = \chi_\lambda,
\]
and hence by (2)
\weq{4}{\sqrt[m]{\psi_\lambda} = \frac{\chi_\lambda}{\left(\sqrt[m]{\psi_1}\right)^{m-\lambda}} = \frac{\chi_\lambda\left(\sqrt[m]{\psi_1}\right)^{\lambda}}{\psi_1},}
and by this, if $\varepsilon$ denotes a fixed primitive $m$th root of unity, all the radicals occurring in (3) are expressed rationally in terms of one of them, $\sqrt[m]{\psi_1}$.

If one gives this one its $m$ different values, one obtains from (3) precisely the $m$ different values $\alpha$.

There is only one exceptional case in which this procedure is not applicable, namely when $\psi_1 = 0$. But by a small modification of the procedure we can free ourselves from such an exceptional case. To this we premise the following.

Let $p$ be a prime dividing $m$, and $m = pn$; as above, let $\varepsilon$ be any fixed primitive $m$th root of unity. Then there is always a $\lambda$ not divisible by $p$ such that $(\varepsilon^{\lambda},\alpha)$ is different from zero. For if we form the difference $\alpha_n - \alpha_0$ according to the formula \S.\,164, (4), we obtain
\[
m(\alpha_n - \alpha_0) = \sum_{0,\,m-1}^{\lambda}(\varepsilon^{-n\lambda} - 1)(\varepsilon^{\lambda},\alpha).
\]
Now, however, $\varepsilon^{-n\lambda} - 1$ is always $= 0$ whenever $\lambda$ is divisible by $p$, and if $(\varepsilon^{\lambda},\alpha)$ vanishes in all the other cases, that is, where $\lambda$ is not divisible by $p$, then $\alpha_n = \alpha_0$, contrary to the hypothesis that the $\alpha$ are all to be different. There is therefore at least one $\lambda$ not divisible by $p$ such that $(\varepsilon^{\lambda},\alpha)$ is different from zero.

We now decompose $m$ into its prime factors and set
\[
m = p_1 p_2 \dots,
\]
%CONT>>
